\hypertarget{the-decision-this-proposal-supports}{%
\section{The decision this proposal
supports}\label{the-decision-this-proposal-supports}}

A universal bank needs to know when a proposed action is likely to
improve a financial outcome, which assumptions make that answer
possible, and when the available information cannot distinguish
competing decisions. The proposed investment is a small reusable
causal-analysis service supporting four journeys: funding responses,
sequential servicing, intervention queries after causal discovery, and
the construction of trustworthy observations. Its immediate deliverable
is a bounded prototype and an evidence-based recommendation about what
to build, reuse or leave unsupported.

The recommendation is to develop and compare a narrow set of these
operations using public records, documented public experiments and
generated cases with known truth. Existing local research supplies
selected working constructions and substantial negative evidence. It
does not yet establish a complete product or a positive return on
deployment. Expensive bank experiments should follow a prototype result
that makes their cost worth considering. The absence of bank data,
institutional sign-off or production access does not prevent completion
of the present proposal. \citep{sp13_scope}

This brief accompanies the 18 October 2026 working monograph, which has
1,060 pages and 751 audited citation records, including 197 genuine
scholarly PDFs. The software, LLM/agent and integration comparisons have
expanded substantially, but their coverage is still incomplete.
Aggregate length does not establish completeness. Literature coverage
retains the 7 October 2026 cutoff, with later acquisitions and uncertain
historical versions explicitly distinguished. The recommendation remains
provisional where evidence or comparison is incomplete.
\citep{sp18_reader}

\hypertarget{the-reusable-product-and-its-customer}{%
\subsection{The reusable product and its
customer}\label{the-reusable-product-and-its-customer}}

The first user is an analyst who can state a banking decision and
inspect its assumptions. The reusable input records a target population,
action and comparator, information available before the action, outcome
and horizon, measurement units, observation process, causal restrictions
and sampling unit. An offer, its receipt and the later balance are
different variables. A measured association is not automatically the
effect of changing an offer.

The output keeps the target attached to its evidence. It may be an
identified formula, an estimate with justified uncertainty, a set of
compatible effects, an informative refusal or a failed computation.
Unresolved graph endpoints, unbounded confidence sets and numerical
failures require different responses. The library, command line and MCP
interface should preserve these distinctions so that a downstream
application can decide whether to act, ask for information or fall back
to its existing workflow. A callable endpoint alone does not establish
identification, useful precision or interoperability.
\citep{sp13_interfaces}

The shared layer would organize three workspaces: intervention studies,
sequential decisions and dynamic systems. The four banking journeys pass
through these workspaces, with observation quality affecting all three.
This organization avoids building a separate estimation stack for every
business team while retaining each team's population, timing and
economic account. The initial service would require analyst review of
consequential scientific choices. No autonomous bank action is proposed.

\hypertarget{a-concrete-first-slice}{%
\subsection{A concrete first slice}\label{a-concrete-first-slice}}

A useful first demonstration begins with a public or generated business
brief, constructs an explicit question and alternative diagrams, checks
the available identification route, and returns a result whose
assumptions and information limits remain visible through the library,
CLI and MCP. For a randomized offer, the first target can be the effect
of assignment. A receipt effect requires additional compliance and
exclusion assumptions. For a confounded observational table, the correct
result may be a documented ambiguity rather than a number.

The comparison must include an expert-designed template and a directly
composed existing-package workflow. The new service earns its place only
if the common representation and checks reduce consequential mistakes or
repeated effort without discarding valid alternatives. A generated
example can expose a wrong answer cheaply; it cannot establish how often
the same failure occurs at a bank.

\hypertarget{what-the-evidence-currently-supports}{%
\section{What the evidence currently
supports}\label{what-the-evidence-currently-supports}}

\hypertarget{funding-and-observations}{%
\subsection{Funding and observations}\label{funding-and-observations}}

Earlier registered synthetic funding work found a passing construction
under a demanding information design: known randomized action moments,
truthful sampled links, known inclusion probabilities and independent
calendars. Its six valid cells passed the prespecified minimum-coverage
screen. False exogeneity, wrong inclusion probabilities and label errors
retained adverse findings. This supports a selected procedure
conditional on its information law, rather than a general correction for
one dependent bank panel. \citep{sp13_funding}

The newer joint-dependence argument makes that limitation more explicit.
A finite confidence-region construction can use declared nuisance bounds
and conservative projection. Its representative precision is inadequate
under the reviewed design. An exact persistence argument also shows why
an unrestricted claim cannot obtain useful shrinking uncertainty merely
by lengthening a single highly persistent series. The practical next
choices are to narrow the dependence family, obtain justified additional
information, or accept wider and possibly unbounded results. The new
derivation has not been turned into a qualified numerical projection
service. \citep{sp13_persistence}

Observation linkage is part of the scientific problem. If different
customer histories are observationally compatible but imply different
intervention values, more computation on the same summaries cannot
choose among them. Public/generated linkage cases can test whether the
prototype preserves this ambiguity. A later bank study would need
representative truth labels and an explicit inclusion mechanism; labels
with unknown selection do not establish the same result.
\citep{sp13_public,sp13_persistence}

\hypertarget{servicing-and-learned-policies}{%
\subsection{Servicing and learned
policies}\label{servicing-and-learned-policies}}

Selected earlier servicing studies support fixed-policy inference in a
restricted nonlinear family under truthful state audits, supplied
historical action probabilities and independent customer trajectories.
Their wrong-logger and information-leakage cases remain failures. They
do not establish the value of an unrestricted learned policy with hidden
borrower state. \citep{sp13_servicing}

A new argument handles estimated action and observation probabilities in
a finite, prespecified history table with bounded rewards. It separates
training from evaluation and makes the support and observation
assumptions explicit. Exact calculations check its policy-value
identities and counterexamples. Its conservative uncertainty can
nevertheless exceed the complete target range at a representative sample
size. The public servicing design uses continuous histories and
unbounded Gaussian/quadratic quantities, so this finite-table theorem
cannot simply be attached to that design. A useful next study needs a
reviewed tighter argument, a separately declared bounded target, or a
different experiment. It must preserve the failed predecessor rather
than change the target after evaluation.
\citep{sp13_servicing_successor}

\hypertarget{discovery-and-agent-support}{%
\subsection{Discovery and agent
support}\label{discovery-and-agent-support}}

Graph discovery constrains possible relationships under stated
assumptions. It may return an equivalence class whose members imply
different interventions. The proposed service must retain that class and
distinguish uncertainty in graph learning from uncertainty in an effect.
Exact latent-risk examples in the current comparison exhibit identical
observed Gaussian distributions with action effects of zero and three
fifths. They prove nonidentification within a class containing both
models; they do not enumerate the entire identified set.
\citep{exec13_latent_examples}

LLMs are proposed as aids for translating business language, finding
missing timing definitions, suggesting diagrams and alternatives, and
preparing auditable tool requests. A model-generated diagram is a
candidate assumption, not observed causal truth. The evaluation must
score false confidence, unsupported orientations, omitted alternatives
and appropriate requests for information. Deterministic task fixtures
and exact examples already support some checks; no live-model
superiority or bank decision gain is established by the new task
designs. A live comparison would require a frozen backend, held-out
tasks, reference adjudication, costs and equally informed baselines.
\citep{sp13_public}

\hypertarget{build-reuse-and-commercial-alternatives}{%
\section{Build, reuse and commercial
alternatives}\label{build-reuse-and-commercial-alternatives}}

The economic case depends on reusing capable existing systems and adding
a scientific connection that users currently have to construct
themselves. Reimplementation is justified only when a measured
constraint or missing capability outweighs the extra verification and
maintenance burden.

\begin{longtable}{@{}p{.19\linewidth}p{.37\linewidth}p{.36\linewidth}@{}}
\toprule
Candidate & Proposed role & Evidence still needed \\
\midrule
\endfirsthead
\toprule
Candidate & Proposed role & Evidence still needed \\
\midrule
\endhead
y0 & Restricted integration of symbolic unconditional identification on a declared DAG/ADMG. Conditional queries use an identified joint law and normalization; the native conditional route remains disabled. & Execute the pinned dependency environment and the exact conversion/failure cases. This is a local integration choice, not runtime admission. \\
DoWhy & Preferred candidate for a supplied-graph identification path, retaining its separation of identification and estimation. & Executable admission is deferred pending dependency/license closure and behavioral checks. No unrestricted estimator or discovery adoption. \\
R and Java candidates & Maintain separate language boundaries and preserve package-specific graph and query semantics. & Execute selected adapters and compare matched scientific operations. A reviewed contract or source algorithm is insufficient evidence of interoperability. \\
Commercial latent engines & Compare supplied-model fitting, measurement, clustering and trial-specific effects with the same configured targets. & Verify edition, interface, rights, fitting diagnostics and actual execution before a licensed trial or procurement recommendation. \\
Shared causal layer & Build the question, assumptions, graph ambiguity, result qualifiers and evidence connections around selected engines. & Demonstrate a complete useful workflow and its failure behavior against expert composition and templates. \\
\bottomrule
\end{longtable}

\citep{sp13_choices,sp13_interfaces}

These distinctions matter commercially. The current primary-document
comparison credits GeNIe with parameter learning on a supplied network;
BayesiaLab with latent representations and segmentation alternatives;
Netica with conditional probability learning and a separately documented
predictive structure path; Stata with supplied latent-class fitting; and
SAS CALIS with specified latent structural models and diagnostics. These
systems have modeling capabilities that the proposed shared layer has
not implemented. Their fitted latent variables are not, by that fact
alone, a discovered causal graph class or a measured confounder.
\citep{exec13_latent_genie_parameters,
exec13_latent_bayesia_data,exec13_latent_netica_algorithms,
exec13_latent_stata_lca,exec13_latent_sas_calis}

Mplus is a particularly relevant causal-effect comparator: its mixture
guide documents complier-average effects in randomized trials. The
prototype cannot claim novelty simply by combining a latent class with a
causal target. Its proposed additional value is preserving which
experimental restrictions, population and financial question justify the
operation, then comparing that whole workflow with the documented
alternative. No commercial engine was run for this comparison.
\citep{exec13_latent_mplus}

The October 13 latent-variable comparison preserved unresolved
edition-specific evidence for several vendors; its omissions are a dated
source-specific result, not a current finding that those products lack
causal capabilities. The expanded commercial chapter now separates the
purchasing unit, supported scientific task, interface and rights,
capacity constraints and full-workflow cost. Its current omissions still
include uninspected panel, synthetic-control, transport and
latent-discovery scope, inaccessible selected policy originals, and
uncertain historical edition bytes. Complete parity, unattended
deployment rights and procurement recommendations require narrower
evidence than a public feature page. No paid engine or matched bank
workflow was executed in that chapter.
\citep{sp13_commercial,sp18_commercial}

Cross-repository reuse also remains conditional. FinMarketing,
MacroFinance, BayesFilter and data/learning repositories are prospective
consumers or component providers under the reviewed architecture. The
current work specifies contracts inside this repository; it does not
establish that sibling repositories have adopted them or that their
owners accept the result semantics. \citep{sp13_interfaces}

\hypertarget{what-the-expanded-interface-work-now-demonstrates}{%
\subsection{What the expanded interface work now
demonstrates}\label{what-the-expanded-interface-work-now-demonstrates}}

The proposal now develops request admission, result interpretation,
publication integrity, retention and capacity as parts of the same
workflow. A local standalone-client check made six actual library and
CLI calls across three requests: a conditional expression, an
unsupported sampling setting and a contradicted premise. Their
scientific meanings agreed. Four effect-set examples were preserved
response replays; they were not four fresh calculations. Five malformed
replies were rejected. This evidence supports the inspected local
examples and their declared schemas. It does not establish a running
notebook kernel, a newly executed MCP integration, another repository's
adoption or production service levels. \citep{sp18_clients}

The retention argument is likewise a design result with a finite
reference. Across ten allowed publication/collection orders, a
stale-snapshot collector loses an artifact needed by a live request in
three orders. A coordinated recheck against current roots avoids that
failure in all ten, with nine successful publications and one safe
rejection. This motivates pins, a live-root check and an explicit rule
for stale publishers. It is not an estimated failure probability or an
implemented storage collector. The prototype should test those semantics
under realistic interruption and restart before an operating guarantee
is proposed. \citep{sp18_retention}

\hypertarget{what-bayesian-advice-adds-and-what-it-cannot-justify}{%
\subsection{What Bayesian advice adds, and what it cannot
justify}\label{what-bayesian-advice-adds-and-what-it-cannot-justify}}

A corrected University of Sao Paulo thesis, published before the cutoff,
supplies a useful adverse comparison for LLM-assisted graph discovery.
Its reported DiBS graph-distance results improve with a GPT prior on
Sachs but worsen on MAGIC-7, MAGIC-ORIGINAL and the authored Summer
case. A DeepSeek prior improves the Summer result while worsening the
two smaller reported alternatives. These are the thesis's tabulated
means, not new experimental runs or significance tests. The integrated
discussion retains its selected feedback experiments, omitted adverse
trials and appendix reproducibility discrepancies. It therefore supports
an experiment comparing information channels, rather than a claim that
LLM advice improves discovery universally. \citep{exec18_videira_thesis}

Repeated model answers estimate response behavior under the query
mechanism. A model can repeat a correct or an incorrect answer equally
consistently. The exact examples in the monograph also show that
restricting an edge model to legal graphs changes its marginal
probabilities, and that a Bayesian feedback update is justified only by
its stipulated likelihood. A proposed deposit-diagram assistant should
preserve response frequencies, semantic categories, admitted graph
support and the evidence for adviser accuracy separately. Adopt these
design lessons now; defer executable DiBS, ancestral-graph or DSPy
integration until pinned code, rights, prompts and experiment artifacts
have been inspected.
\citep{exec18_bayesian_reference,exec18_videira_thesis}

\hypertarget{public-evidence-before-expensive-bank-experiments}{%
\section{Public evidence before expensive bank
experiments}\label{public-evidence-before-expensive-bank-experiments}}

Six reviewed study designs give the prototype inexpensive routes
forward. They are designs with local checks, not newly executed or
prospectively registered final studies. Before final evaluation, each
selected route must freeze its code and input identities, targets,
splits or seeds, resource budget, acceptance rule and failure
accounting. Optional dataset access cannot block an independent
generated route. \citep{sp13_public}

\begin{longtable}{@{}p{.22\linewidth}p{.36\linewidth}p{.34\linewidth}@{}}
\toprule
Question & Public or generated evidence & What it can establish \\
\midrule
\endfirsthead
\toprule
Question & Public or generated evidence & What it can establish \\
\midrule
\endhead
Formulation and timing & Archived 401(k) eligibility/participation and UCI bank-marketing task briefs, with explicit pre-action fields. & Preservation of target, clock and alternatives. A predictive marketing table is not treatment-effect ground truth. \\
Funding and panels & A bounded FDIC vintage with verified dictionary and unit map; an independent local macro/generated route. & Reproducible units and conditional/descriptive analyses. Extraction and vintage questions remain local prerequisites. \\
Single-decision effects & A generated randomized law with arm means, interacted regression and cross-fitted augmented weighting. & Matched estimation and uncertainty behavior in the declared law; retain overlap and misspecification failures. \\
Sequential servicing & Controlled histories, actions, observation loss and fixed versus trained policies. & Value-estimation and clock diagnostics under a reviewed law. The continuous design needs its own uncertainty argument. \\
Graphs and linkage & Generated graph-equivalence and entity-corruption cases with evaluator-held truth. & Detection of invented identification and preservation of uncertainty; no inference of bank label quality. \\
Agent workflow & Frozen public/generated briefs, templates, task-family splits and an ambiguity rubric. & Offline task and interface correctness. Live-model performance requires an actual matched held-out run. \\
\bottomrule
\end{longtable}

The suggested first empirical slice pairs formulation/timing cases with
the generated randomized-effect comparison. It is sufficiently narrow to
compare complete workflows and has a known-truth route without paid
data. Add graph and linkage cases as deliberate refusal/ambiguity tests.
Servicing and dependent funding arguments can advance independently, but
their unresolved precision problems must not be hidden by a successful
simpler demonstration.

For every offered task, record the scientific result, harmful confident
errors, useful refusals, analyst corrections, time and resource
consumption. Compare the same information, target and failure policy.
Retain failed or unavailable requests in the denominator. A successful
average can otherwise conceal an unsafe or scientifically misleading
subgroup. No elapsed-time improvement is automatically a measured labor
saving or a causal financial gain. \citep{sp13_economics}

\hypertarget{economics-and-the-next-experiment}{%
\section{Economics and the next
experiment}\label{economics-and-the-next-experiment}}

Let \(N\) be the number of attempted decision tasks in one year, \(D\)
the incremental gross monetary utility of the proposed action relative
to the chosen baseline, \(F\) fixed incremental cost, \(c\) variable
incremental cost per attempted task, and \(R\) an additional risk cost
not already counted in \(D\). A transparent planning account is \[
 V=N\operatorname{E}[D]-F-Nc-R.
\] The baseline must be a feasible alternative under the same
information. A correct refusal may leave the baseline action unchanged
while still costing review time. A misleading recommendation can have
negative utility. Loss must not be counted in both \(D\) and \(R\), and
unacceptable actions remain excluded by constraints irrespective of
their average financial payoff. \citep{sp13_economics}

\hypertarget{a-reproducible-illustration-with-assumed-inputs}{%
\subsection{A reproducible illustration, with assumed
inputs}\label{a-reproducible-illustration-with-assumed-inputs}}

The following one-year illustration uses USD solely as a declared
planning unit. Assume 10,000 attempted tasks, fixed incremental cost of
USD 100,000, variable incremental cost of USD 6 per task and no separate
\(R\). All volumes, prices, gains, probabilities and signal likelihoods
below are assumptions; none is a measured bank benefit or vendor quote.
The calculation excludes discounting because its horizon is one year.
\citep{sp13_economics}

\begin{center}
\begin{tabular}{@{}lrrrr@{}}
\toprule
Regime & Prior & Gross gain/task & Net value & Pr(positive signal) \\
\midrule
High & 0.4 & USD 40 & USD 240,000 & 0.8 \\
Moderate & 0.4 & USD 20 & USD 40,000 & 0.5 \\
Adverse & 0.2 & USD $-2$ & USD $-180,000$ & 0.1 \\
\bottomrule
\end{tabular}
\end{center}

Expected incremental value is USD 76,000 before further information.
This number alone cannot justify a real investment because its inputs
are assumed. Under the same assumptions, an experiment produces a
positive signal with probability 0.54. Expected net value conditional on
that signal is approximately USD 150,370, and conditional on a negative
signal it is approximately minus USD 11,304. Choosing to build only
after a positive signal gives expected value USD 81,200 before the
experiment's own cost.

The expected value of sample information is consequently USD 5,200. An
experiment costing USD 8,000 loses USD 2,800 in expected information
value; at USD 3,000 it gains USD 2,200, assuming identical signal
informativeness. Perfect information has value USD 36,000 in this
illustration. That is an upper bound on the information component of
spending for this particular decision and prior, not a universal
experiment budget. The prototype must help assess the assumed signal
quality before these calculations can inform an expensive study.
\citep{sp13_economics}

Holding other assumptions fixed, break-even gross gain is USD 16 per
attempted task. Halving volume raises it to USD 26; doubling volume
reduces it to USD 11. An additional USD 4 of review cost per task raises
the 10,000-task threshold to USD 20. These sensitivities are arithmetic
consequences, not estimates of demand or analyst effort. Refusals and
failed requests still incur applicable costs and remain among the
attempted tasks. \citep{sp13_economics_sensitivity}

\hypertarget{costs-to-measure-in-the-prototype}{%
\subsection{Costs to measure in the
prototype}\label{costs-to-measure-in-the-prototype}}

The measured account should separate data acquisition and cleaning,
labels and identity audits, model and tool calls, analyst formulation
and checking, package setup, integration and maintenance, compute and
storage, incident recovery and independent validation. Each cost needs
its unit, quantity, date/version, price source and incremental share
relative to the baseline. The current actual-price fields remain
unfilled. Public examples can reveal setup burden and local resource use
without establishing bank-scale service levels, customer outcomes or
willingness to adopt. \citep{sp13_economics}

A later bank experiment should state the decisions it distinguishes, the
uncertain assumption that matters, the target population and horizon,
expected signals, losses from wrong decisions, and the full cost of
obtaining the information. It should compare cheaper public/generated
alternatives first. If no feasible experiment has plausible positive net
information value, narrow or defer the capability. A demand for
authentic evidence becomes appropriate when that specific experiment is
justified; it is not a prerequisite for writing the current proposal.

\hypertarget{what-remains-before-a-complete-proposal}{%
\section{What remains before a complete
proposal}\label{what-remains-before-a-complete-proposal}}

The software comparison still has 80 open and 20 inconclusive cells in
its 102-cell coverage matrix. The reuse register has 99 open and one
inconclusive decision among 115 candidates. These counts describe
unresolved review obligations, not numbers of defective products. Each
closure needs an explicit original, supported scope, comparison and
adopt/integrate/port/reject decision. No absence finding follows from an
unread source. \citep{sp18_progress}

Twenty of 22 chapter reviews meet their current author criteria. L1 and
L6 remain inconclusive because frontier and framework coverage is
incomplete despite meeting minimum depth. The original-candidate census
now contains 40 inspected and 73 unread originals out of 113;
citation-chain leads retain their own denominator. Fourteen task reviews
and all twelve integration requirements have scoped proposal outcomes,
but those receipts do not establish a live model benchmark or
cross-repository adoption. The three mandatory repairs remain open. The
agreed depth contracts retain 100--150 substantive software pages,
approximately 100 agent pages and a developed integration treatment;
repeated material does not earn duplicate credit.
\citep{sp18_progress,sp18_census,sp13_scope}

The current sponsor update reconciles the recommendation, selected new
evidence and remaining work to this reader. Final proposal acceptance
still requires complete coverage reconciliation and rendered visual
review of the exact release. Source audits, exact arithmetic and
all-page text bounds are narrower checks and cannot substitute for
visual acceptance. Missing bank data, platform access or signed adoption
remain later conditional needs; they do not stop independent proposal
research. \citep{sp18_reader,sp18_progress}

The evidence routes are also different. We know how to complete missing
source comparisons, write and test selected adapters, run registered
generated benchmarks and measure local workflow costs. Their success
still requires execution and review. We do not presently have a general
useful solution for arbitrary latent graphs, unrestricted persistent
dependence, hidden-state optimal policies or arbitrary nuisance
learners. The proposal can be complete with explicit bounded claims and
negative findings in these areas; it cannot claim their scientific
resolution. \citep{sp18_reader,
sp13_persistence,sp13_servicing_successor}

The present recommendation is therefore to test the narrow common
service against existing composition, use that evidence to refine the
investment scenario, and request a bank experiment only for a
consequential question that the cheaper evidence cannot settle.
Independent scientific acceptance, target-bank validation, production
reliability and adoption are later decisions with their own evidence.
They remain distinct from completing a credible proposal and a useful
bounded prototype.
